\documentclass[11pt]{article}
\usepackage[margin=1in]{geometry}
\usepackage{amsmath, amssymb, amsthm}
\usepackage{graphicx, booktabs}
\usepackage[shortlabels]{enumitem}
\usepackage{hyperref}
\usepackage{fancyhdr}
\usepackage{microtype}
\usepackage{tikz}
\usetikzlibrary{arrows.meta, intersections}

\newtheorem{theorem}{Theorem}[section]
\newtheorem{proposition}[theorem]{Proposition}
\theoremstyle{definition}
\newtheorem{definition}[theorem]{Definition}
\newtheorem{example}[theorem]{Example}
\theoremstyle{remark}
\newtheorem*{remark}{Remark}

\newcommand{\MC}{\mathrm{MC}}
\newcommand{\MR}{\mathrm{MR}}
\newcommand{\CS}{\mathrm{CS}}
\newcommand{\PS}{\mathrm{PS}}
\newcommand{\Qd}{Q^{d}}
\newcommand{\Qs}{Q^{s}}

\pagestyle{fancy}
\fancyhf{}
\lhead{ECON 101 --- Introduction to Economics}
\chead{Lecture 2}
\rhead{Fall 2026}
\cfoot{\thepage}

% -------------------------------------------------------
% Helper: basic S&D diagram
% Usage: \SDdiagram{title}{extra tikz code}
\newcommand{\blankSDdiagram}[1]{%
\begin{center}
\begin{tikzpicture}[scale=0.9,
  >=Stealth,
  every node/.style={font=\small}]
  % Axes
  \draw[->] (0,0) -- (6.5,0) node[right] {$Q$};
  \draw[->] (0,0) -- (0,6.5) node[above] {$P$};
  #1
\end{tikzpicture}
\end{center}}

\begin{document}

\begin{center}
  {\LARGE\bfseries ECON 101: Introduction to Economics}\\[0.4em]
  {\large Lecture 2: Supply and Demand}
\end{center}
\medskip\hrule\medskip

% -------------------------------------------------------
\section{Motivation}

Last lecture we established that economics is about allocating scarce resources.
\textit{How} does a decentralised economy coordinate those decisions?
The answer, in most markets, is the \textbf{price system}: prices rise and fall
until the quantity sellers want to sell equals the quantity buyers want to buy.
Supply and demand is the workhorse model for understanding that process.

% -------------------------------------------------------
\section{Demand}

\begin{definition}
  The \textbf{quantity demanded} of a good is the amount buyers are willing and
  able to purchase at a given price, holding all other factors constant
  (\textit{ceteris paribus}).
\end{definition}

\begin{definition}
  The \textbf{Law of Demand}: as the price of a good rises, the quantity demanded
  falls; as the price falls, the quantity demanded rises.  All else equal,
  $P \uparrow \Rightarrow \Qd \downarrow$.
\end{definition}

Two complementary reasons underlie the law of demand:
\begin{itemize}
  \item \textbf{Substitution effect:} a higher price makes other goods relatively
    cheaper, so buyers substitute away.
  \item \textbf{Income effect:} a higher price erodes real purchasing power,
    reducing consumption of normal goods.
\end{itemize}

\subsection{The Demand Curve}

A \textbf{demand curve} plots the quantity demanded at each price.  It slopes
downward by the law of demand.  A linear example:
\[
  \Qd = a - bP, \qquad a,b > 0.
\]

\blankSDdiagram{%
  % Demand line: passes through (1,5) and (5,1)
  \draw[blue, thick, name path=D] (0.5,5.5) -- (5.5,0.5)
    node[right] {$D$};
  \node[blue] at (1.2,5.2) {};
}

\subsection{Shifters of Demand}

A \textit{change in price} causes a \textbf{movement along} the demand curve.
A \textit{change in any other determinant} causes a \textbf{shift} of the entire
curve.  The main shifters are:

\begin{enumerate}[(1)]
  \item \textbf{Income.} For a \textit{normal good}, higher income $\Rightarrow$
    demand shifts right.  For an \textit{inferior good}, it shifts left.
  \item \textbf{Prices of related goods.}
    \begin{itemize}
      \item \textit{Substitutes} (e.g.\ coffee and tea): $P_{\text{coffee}} \uparrow
        \Rightarrow D_{\text{tea}}$ shifts right.
      \item \textit{Complements} (e.g.\ coffee and cream): $P_{\text{coffee}} \uparrow
        \Rightarrow D_{\text{cream}}$ shifts left.
    \end{itemize}
  \item \textbf{Tastes and preferences.}  Advertising, trends, new information.
  \item \textbf{Expectations.}  If buyers expect a future price rise, current
    demand increases.
  \item \textbf{Number of buyers.}  More buyers $\Rightarrow$ market demand shifts
    right.
\end{enumerate}

\begin{example}
  A health report links red meat to cardiovascular disease.  Tastes shift: demand
  for beef falls (curve shifts left) and demand for chicken rises (curve shifts
  right), even though neither price changed.
\end{example}

% -------------------------------------------------------
\section{Supply}

\begin{definition}
  The \textbf{quantity supplied} is the amount sellers are willing and able to
  offer for sale at a given price, ceteris paribus.
\end{definition}

\begin{definition}
  The \textbf{Law of Supply}: as price rises, quantity supplied rises; as price
  falls, quantity supplied falls.  $P \uparrow \Rightarrow \Qs \uparrow$.
\end{definition}

The intuition: a higher price raises the profitability of production, inducing
existing firms to expand output and attracting new entrants.

\subsection{The Supply Curve}

A \textbf{supply curve} slopes upward.  A linear example:
\[
  \Qs = -c + dP, \qquad c,d > 0, \quad P \geq c/d.
\]

\blankSDdiagram{%
  % Supply line: passes through (1,1) and (5,5)
  \draw[red, thick, name path=S] (0.5,0.5) -- (5.5,5.5)
    node[right] {$S$};
}

\subsection{Shifters of Supply}

\begin{enumerate}[(1)]
  \item \textbf{Input prices.}  Cheaper inputs $\Rightarrow$ supply shifts right.
  \item \textbf{Technology.}  Improved technology lowers production costs and
    shifts supply right.
  \item \textbf{Expectations.}  Sellers who expect higher future prices may
    withhold supply today (shift left now).
  \item \textbf{Number of sellers.}  More producers $\Rightarrow$ market supply
    shifts right.
  \item \textbf{Government policies.}  Taxes raise costs (shift left); subsidies
    lower costs (shift right).
\end{enumerate}

% -------------------------------------------------------
\section{Market Equilibrium}

\begin{definition}
  A market is in \textbf{equilibrium} when quantity demanded equals quantity
  supplied:
  \[
    \Qd(P^*) = \Qs(P^*).
  \]
  The price $P^*$ is the \textbf{equilibrium price} (market-clearing price) and
  $Q^* \equiv \Qd(P^*)$ is the \textbf{equilibrium quantity}.
\end{definition}

\subsection{Solving for Equilibrium}

Given
\[
  \Qd = a - bP \quad \text{and} \quad \Qs = -c + dP,
\]
set $\Qd = \Qs$:
\[
  a - bP^* = -c + dP^* \implies P^* = \frac{a+c}{b+d}, \qquad
  Q^* = a - b\cdot\frac{a+c}{b+d} = \frac{ad - bc}{b+d}.
\]

\blankSDdiagram{%
  % Demand
  \draw[blue, thick, name path=D] (0.5,5.5) -- (5.5,0.5)
    node[right] {$D$};
  % Supply
  \draw[red, thick, name path=S] (0.5,0.5) -- (5.5,5.5)
    node[right] {$S$};
  % Equilibrium intersection
  \fill (3,3) circle (2.5pt) node[above right] {$(Q^*, P^*)$};
  % Dashed lines to axes
  \draw[dashed, gray] (3,3) -- (3,0) node[below] {$Q^*$};
  \draw[dashed, gray] (3,3) -- (0,3) node[left]  {$P^*$};
}

\subsection{Surplus and Shortage}

\begin{itemize}
  \item \textbf{Surplus} ($P > P^*$): $\Qs > \Qd$.  Unsold inventory
    accumulates; sellers cut prices, restoring equilibrium.
  \item \textbf{Shortage} ($P < P^*$): $\Qd > \Qs$.  Buyers bid prices up,
    restoring equilibrium.
\end{itemize}

The market mechanism self-corrects: deviations from $P^*$ generate pressures
that push prices back toward equilibrium.

% -------------------------------------------------------
\section{Comparative Statics}

Comparative statics asks: \textit{if one determinant changes, how do $P^*$ and
$Q^*$ change?}  The procedure is:
\begin{enumerate}[(1)]
  \item Determine whether the event shifts $D$, $S$, or both, and in which
    direction.
  \item Draw the new curve.
  \item Read off the change in equilibrium.
\end{enumerate}

\begin{example}[Demand increase]
  A rise in consumer income (normal good) shifts $D$ rightward to $D'$,
  raising both $P^*$ and $Q^*$.

  \blankSDdiagram{%
    % Supply
    \draw[red, thick] (0.5,0.5) -- (5.5,5.5) node[right] {$S$};
    % Original demand
    \draw[blue, thick] (0.5,5.5) -- (5.5,0.5) node[right] {$D$};
    \fill (3,3) circle (2pt);
    \draw[dashed, gray] (3,3) -- (3,0) node[below] {$Q_1^*$};
    \draw[dashed, gray] (3,3) -- (0,3) node[left]  {$P_1^*$};
    % New demand (shifted right by 1.5 units)
    \draw[blue, thick, dashed] (2,5.5) -- (6.2,1.3) node[right] {$D'$};
    \fill (4.25,4.25) circle (2pt);
    \draw[dashed, gray] (4.25,4.25) -- (4.25,0) node[below] {$Q_2^*$};
    \draw[dashed, gray] (4.25,4.25) -- (0,4.25) node[left] {$P_2^*$};
  }
\end{example}

\begin{example}[Supply decrease]
  An increase in input costs shifts $S$ leftward to $S'$, raising $P^*$ but
  lowering $Q^*$.

  \blankSDdiagram{%
    % Demand
    \draw[blue, thick] (0.5,5.5) -- (5.5,0.5) node[right] {$D$};
    % Original supply
    \draw[red, thick] (0.5,0.5) -- (5.5,5.5) node[right] {$S$};
    \fill (3,3) circle (2pt);
    \draw[dashed, gray] (3,3) -- (3,0) node[below] {$Q_1^*$};
    \draw[dashed, gray] (3,3) -- (0,3) node[left]  {$P_1^*$};
    % New supply (shifted left by 1.5 units)
    \draw[red, thick, dashed] (2,0.5) -- (5.5,4) node[right] {$S'$};
    \fill (2.35,4.15) circle (2pt);
    \draw[dashed, gray] (2.35,4.15) -- (2.35,0) node[below] {$Q_2^*$};
    \draw[dashed, gray] (2.35,4.15) -- (0,4.15) node[left] {$P_2^*$};
  }
\end{example}

Table~\ref{tab:comstats} summarises all four single-shift cases.

\begin{table}[h]
\centering
\caption{Comparative statics: direction of change in equilibrium}
\label{tab:comstats}
\begin{tabular}{lcc}
\toprule
Event & $\Delta P^*$ & $\Delta Q^*$ \\
\midrule
Demand increases ($D$ shifts right) & $+$ & $+$ \\
Demand decreases ($D$ shifts left)  & $-$ & $-$ \\
Supply increases ($S$ shifts right) & $-$ & $+$ \\
Supply decreases ($S$ shifts left)  & $+$ & $-$ \\
\bottomrule
\end{tabular}
\end{table}

\begin{remark}
  When \textit{both} curves shift simultaneously, the direction of one of the
  two equilibrium variables becomes ambiguous without knowing the magnitudes of
  the shifts.
\end{remark}

% -------------------------------------------------------
\section{Applications: Price Controls}

Governments sometimes prevent prices from reaching $P^*$.

\begin{definition}
  A \textbf{price ceiling} is a legal maximum price.  It is binding only when
  set below $P^*$, creating a \textbf{shortage}: $\Qd > \Qs$.

  A \textbf{price floor} is a legal minimum price.  It is binding only when
  set above $P^*$, creating a \textbf{surplus}: $\Qs > \Qd$.
\end{definition}

\begin{example}[Rent control]
  A city caps apartment rent below the market-clearing level (a binding price
  ceiling).  Quantity demanded of apartments exceeds quantity supplied, producing
  a housing shortage.  Non-price rationing mechanisms (waiting lists, connections)
  emerge.
\end{example}

\begin{example}[Minimum wage]
  A binding minimum wage sets the wage floor above the market wage.
  Quantity of labour supplied exceeds quantity demanded, producing unemployment.
\end{example}

% -------------------------------------------------------
\section{Summary}

\begin{itemize}
  \item The \textbf{law of demand}: $P \uparrow \Rightarrow \Qd \downarrow$
    (downward-sloping demand curve).
  \item The \textbf{law of supply}: $P \uparrow \Rightarrow \Qs \uparrow$
    (upward-sloping supply curve).
  \item \textbf{Equilibrium} ($\Qd = \Qs$) determines $P^*$ and $Q^*$.
    Surpluses and shortages push prices back to $P^*$.
  \item Shifts in $D$ or $S$ change equilibrium; use comparative statics to
    determine the direction.
  \item \textbf{Price controls} prevent market clearing, generating persistent
    shortages (ceilings) or surpluses (floors).
\end{itemize}

\end{document}
